\documentclass[a4paper,12pt]{article}

\usepackage[T2A]{fontenc}
\usepackage[utf8]{inputenc}
\usepackage[russian]{babel}
\usepackage{amsmath,amssymb,amsthm,mathtools}
\usepackage{helvet}
\renewcommand{\familydefault}{\sfdefault}
\usepackage{icomma}
\usepackage[a4paper,margin=2.1cm]{geometry}
\usepackage{microtype}
\usepackage{tikz}
\usetikzlibrary{calc}
\usepackage{tkz-euclide}
\usepackage{pgfplots}
\pgfplotsset{compat=1.18}
\usepackage{tabularx,booktabs,longtable}
\usepackage{enumitem}
\usepackage{hyperref}
\usepackage{parskip}
\usepackage{fancyhdr}
\usepackage{setspace}
\usepackage{xcolor}

\definecolor{accent}{HTML}{008080}
\definecolor{darkaccent}{HTML}{004D4D}
\definecolor{textgray}{HTML}{333333}
\definecolor{softgray}{HTML}{F3F5F5}

\hypersetup{
  colorlinks=true,
  linkcolor=darkaccent,
  urlcolor=darkaccent,
  pdftitle={Чек-лист: показательные уравнения},
  pdfauthor={Лёвин Артём Александрович}
}

\setlength{\headheight}{25pt}
\onehalfspacing
\setlength{\parindent}{0pt}
\setlength{\parskip}{0.55em}
\setlist[itemize]{leftmargin=1.5em,itemsep=0.25em,topsep=0.25em}
\setlist[enumerate]{leftmargin=1.8em,itemsep=0.35em,topsep=0.25em}
\renewcommand{\arraystretch}{1.25}

\pagestyle{fancy}
\fancyhf{}
\fancyhead[L]{\small\textcolor{textgray}{Показательные уравнения}}
\fancyhead[R]{\small\textcolor{textgray}{Анастасия Павлова}}
\fancyfoot[C]{\small\thepage}
\renewcommand{\headrulewidth}{0.4pt}
\renewcommand{\headrule}{\hbox to\headwidth{\color{accent}\leaders\hrule height \headrulewidth\hfill}}

\newcommand{\sectionline}{\par\vspace{0.15em}\noindent\textcolor{accent}{\rule{\linewidth}{0.7pt}}\par\vspace{0.2em}}
\newcommand{\key}[1]{\textcolor{darkaccent}{\textbf{#1}}}

\begin{document}

% =========================================================
% ТИТУЛЬНЫЙ ЛИСТ
% =========================================================
\begin{titlepage}
\thispagestyle{empty}
\centering
\vspace*{2.5cm}

{\Large\textcolor{darkaccent}{ЧЕК-ЛИСТ}\par}
\vspace{0.7cm}

{\Huge\bfseries Показательные уравнения\par}

\vspace{0.35cm}

{\Large Приведение к общему основанию, свойства степеней,\par}
{\Large дроби, корни и вынесение общего множителя\par}

\vspace{1.3cm}

\textcolor{accent}{\rule{0.72\linewidth}{1pt}}

\vspace{1.3cm}

{\large ЕГЭ по профильной математике\par}

\vspace{0.5cm}

{\large Лёвин Артём Александрович\par}
{\large эксклюзивно для Анастасии Павловой\par}

\vspace{0.7cm}

{\large \today\par}

\vfill

\begin{tikzpicture}[remember picture,overlay]
  \draw[accent,line width=1.1pt]
    ($(current page.south west)+(1.8cm,1.55cm)$)
      .. controls ($(current page.south west)+(6.2cm,2.45cm)$)
      and ($(current page.south east)+(-6.2cm,0.75cm)$)
      .. ($(current page.south east)+(-1.8cm,1.55cm)$);

  \draw[darkaccent,line width=0.45pt]
    ($(current page.south west)+(1.8cm,1.28cm)$)
      .. controls ($(current page.south west)+(6.6cm,2.05cm)$)
      and ($(current page.south east)+(-6.6cm,0.88cm)$)
      .. ($(current page.south east)+(-1.8cm,1.28cm)$);
\end{tikzpicture}

\end{titlepage}

% =========================================================
% ОСНОВНОЙ ЧЕК-ЛИСТ
% =========================================================

\section*{1. Главная идея показательного уравнения}
\sectionline

Базовая цель --- привести обе части уравнения к
\key{одинаковому основанию}:

\[
a^{f(x)}=a^{g(x)},
\qquad
a>0,
\quad
a\ne1.
\]

При этих условиях показательная функция взаимно однозначна, поэтому

\[
a^{f(x)}=a^{g(x)}
\quad\Longleftrightarrow\quad
f(x)=g(x).
\]

\key{Важно.}
Переход к равенству показателей допустим только тогда, когда слева и справа осталось по одной степени с одинаковым основанием. Дополнительные множители, делители, слагаемые и вычитаемые предварительно нужно преобразовать.

\subsection*{Пример 1. Простое приведение к основанию}

\[
2^{4-2x}=64.
\]

Так как

\[
64=2^6,
\]

получаем

\[
2^{4-2x}=2^6,
\]

следовательно,

\[
4-2x=6,
\]

\[
-2x=2,
\]

\[
\boxed{x=-1}.
\]

\subsection*{Пример 2. Основание 3}

\[
3^{2x-1}=9.
\]

Представим

\[
9=3^2.
\]

Тогда

\[
3^{2x-1}=3^2,
\]

\[
2x-1=2,
\]

\[
2x=3,
\]

\[
\boxed{x=\frac32}.
\]

% =========================================================

\section*{2. Формулы степеней}
\sectionline

Для положительного основания \(a\) используются преобразования

\[
a^m\cdot a^n=a^{m+n},
\]

\[
\frac{a^m}{a^n}=a^{m-n},
\]

\[
\left(a^m\right)^n=a^{mn},
\]

\[
a^{-n}=\frac1{a^n},
\]

\[
\frac1{a^n}=a^{-n},
\]

\[
\sqrt[k]{a^m}=a^{m/k}.
\]

\key{Главное различие операций:}

\[
a^m\cdot a^n=a^{m+n},
\]

но формулы вида

\[
a^m+a^n=a^{m+n}
\]

или

\[
a^m-a^n=a^{m-n}
\]

не существует.

Сложение и вычитание требуют других преобразований.

% =========================================================

\section*{3. Составные числа}
\sectionline

При решении сразу ищите числа, которые можно представить степенью выбранного основания:

\[
4=2^2,
\qquad
8=2^3,
\qquad
16=2^4,
\qquad
32=2^5,
\qquad
64=2^6,
\]

\[
9=3^2,
\qquad
27=3^3,
\qquad
81=3^4,
\]

\[
25=5^2,
\qquad
125=5^3.
\]

Цель --- получить одинаковые основания в обеих частях уравнения.

% =========================================================

\section*{4. Дроби как степени}
\sectionline

Дробь удобно переводить в отрицательную степень:

\[
\frac1{27}
=
\frac1{3^3}
=
3^{-3}.
\]

Аналогично,

\[
\frac19=3^{-2},
\qquad
\frac1{25}=5^{-2},
\qquad
\frac1{16}=2^{-4}.
\]

Общее правило:

\[
\frac1{a^n}=a^{-n}.
\]

\subsection*{Можно сохранить дробное основание}

Рассмотрим

\[
\left(\frac13\right)^{x-8}=\frac19.
\]

Так как

\[
\frac19=\left(\frac13\right)^2,
\]

можно сразу получить

\[
x-8=2,
\]

\[
\boxed{x=10}.
\]

Можно решить то же уравнение через основание \(3\):

\[
3^{-(x-8)}=3^{-2}.
\]

Оба способа корректны.

\key{Практический принцип:}
если дробное основание уже одинаково в обеих частях, его можно сохранить.

% =========================================================

\section*{5. Когда дробное основание приходится менять}
\sectionline

Рассмотрим

\[
\left(\frac12\right)^{6-2x}=4.
\]

Справа находится целая степень двойки:

\[
4=2^2.
\]

Слева

\[
\frac12=2^{-1},
\]

поэтому

\[
\left(2^{-1}\right)^{6-2x}=2^2.
\]

Перемножаем показатели:

\[
2^{-(6-2x)}=2^2,
\]

\[
-6+2x=2,
\]

\[
2x=8,
\]

\[
\boxed{x=4}.
\]

\key{Идея:}
если основания имеют разные формы, приведите их к одному базовому числу.

% =========================================================

\section*{6. Корни как дробные степени}
\sectionline

Корень переводится в степень по формуле

\[
\sqrt[k]{a^m}=a^{m/k}.
\]

Например,

\[
\sqrt{27}
=
\sqrt{3^3}
=
3^{3/2},
\]

\[
\sqrt[3]{25}
=
\sqrt[3]{5^2}
=
5^{2/3}.
\]

\subsection*{Пример}

\[
3^{2x-1}=\sqrt{27}.
\]

Преобразуем правую часть:

\[
\sqrt{27}=3^{3/2}.
\]

Получаем

\[
3^{2x-1}=3^{3/2}.
\]

Следовательно,

\[
2x-1=\frac32,
\]

\[
2x=\frac52,
\]

\[
\boxed{x=\frac54}.
\]

% =========================================================

\section*{7. Универсальный алгоритм}
\sectionline

\begin{enumerate}
  \item Осмотрите уравнение целиком.
  \item Найдите составные числа, дроби, корни и степени степеней.
  \item Выберите удобное общее основание.
  \item Преобразуйте левую и правую части по отдельности.
  \item Используйте формулы степеней.
  \item Добейтесь вида
  \[
  a^{f(x)}=a^{g(x)}.
  \]
  \item Проверьте условия
  \[
  a>0,
  \qquad
  a\ne1.
  \]
  \item Приравняйте показатели:
  \[
  f(x)=g(x).
  \]
  \item Решите полученное алгебраическое уравнение.
  \item Выполните быструю проверку вычислений.
\end{enumerate}

\key{Полезный приём:}
сложное выражение лучше преобразовывать по небольшим фрагментам, а затем собирать результат.

% =========================================================

\section*{8. Дополнительный множитель перед степенью}
\sectionline

Рассмотрим уравнение

\[
3\cdot9^{x-1}=\frac1{27}.
\]

Если сразу заменить

\[
9=3^2,
\qquad
\frac1{27}=3^{-3},
\]

получим

\[
3\cdot\left(3^2\right)^{x-1}=3^{-3}.
\]

После преобразования степени степени:

\[
3^1\cdot3^{2x-2}=3^{-3}.
\]

Слева находится произведение степеней с одинаковым основанием:

\[
3^{1+2x-2}=3^{-3},
\]

\[
3^{2x-1}=3^{-3}.
\]

Теперь основания можно отбросить:

\[
2x-1=-3,
\]

\[
2x=-2,
\]

\[
\boxed{x=-1}.
\]

\key{Контрольный вопрос:}
осталась ли в каждой части ровно одна степень?

Если перед степенью остался отдельный множитель, преобразование ещё не завершено.

% =========================================================

\section*{9. Сумма и разность степеней}
\sectionline

Если между степенными выражениями стоит \(+\) или \(-\), непосредственно складывать или вычитать показатели нельзя.

Часто нужно привести выражения к одинаковому степенному фрагменту и вынести его как общий множитель.

\subsection*{Пример}

\[
6^{x+1}-6^x=180.
\]

Используем

\[
6^{x+1}=6^x\cdot6.
\]

Тогда

\[
6^x\cdot6-6^x=180.
\]

Выносим общий множитель:

\[
6^x(6-1)=180.
\]

Получаем

\[
5\cdot6^x=180,
\]

\[
6^x=36.
\]

Так как

\[
36=6^2,
\]

то

\[
6^x=6^2,
\]

следовательно,

\[
\boxed{x=2}.
\]

Полезный шаблон:

\[
a^{x+r}\pm a^x
=
a^x\left(a^r\pm1\right).
\]

% =========================================================

\section*{10. Сложные выражения: работа по действиям}
\sectionline

Если одновременно присутствуют корни, дроби, составные основания и сложные показатели, выполняйте преобразования последовательно.

Например,

\[
\sqrt[3]{9}
=
\sqrt[3]{3^2}
=
3^{2/3},
\]

а

\[
27^{\frac{x}{3}+1}
=
\left(3^3\right)^{\frac{x}{3}+1}.
\]

По формуле степени степени:

\[
\left(3^3\right)^{\frac{x}{3}+1}
=
3^{3\left(\frac{x}{3}+1\right)}
=
3^{x+3}.
\]

\key{Следите за скобками.}

Если показатель содержит сумму или разность, весь показатель должен участвовать в умножении:

\[
3\left(\frac{x}{3}+1\right)=x+3.
\]

Удобная последовательность:

\begin{enumerate}
  \item преобразовать корни;
  \item преобразовать дроби;
  \item заменить составные основания;
  \item выполнить степень степени;
  \item упростить показатели;
  \item собрать уравнение.
\end{enumerate}

% =========================================================

\section*{11. Разбор важной ошибки из домашней работы}
\sectionline

Рассмотрим

\[
\frac{6x^{12}+9x^{12}}{9x^{12}}.
\]

Так как в знаменателе присутствует \(x^{12}\),

\[
x\ne0.
\]

\subsection*{Правило сокращения}

Сокращать можно \key{общие множители} числителя и знаменателя.

Нельзя сокращать отдельное слагаемое числителя через знак \(+\) или \(-\).

Например, нельзя сократить только

\[
9x^{12}
\]

из числителя с

\[
9x^{12}
\]

в знаменателе, оставив первое слагаемое без деления.

\subsection*{Путь 1. Сначала привести подобные}

\[
6x^{12}+9x^{12}=15x^{12}.
\]

Тогда

\[
\frac{6x^{12}+9x^{12}}{9x^{12}}
=
\frac{15x^{12}}{9x^{12}}.
\]

При \(x\ne0\):

\[
\frac{15x^{12}}{9x^{12}}
=
\frac{15}{9}
=
\boxed{\frac53}.
\]

\subsection*{Путь 2. Разделить каждый член числителя}

Используем

\[
\frac{A+B}{C}
=
\frac AC+\frac BC.
\]

Получаем

\[
\frac{6x^{12}}{9x^{12}}
+
\frac{9x^{12}}{9x^{12}}
=
\frac23+1
=
\boxed{\frac53},
\qquad
x\ne0.
\]

Первый способ здесь предпочтительнее: выражение быстрее становится компактным.

% =========================================================

\section*{12. Подобные слагаемые}
\sectionline

Подобные слагаемые имеют одинаковую буквенную часть с одинаковыми показателями степеней.

Например,

\[
5t^3+7t^3
\]

--- подобные слагаемые.

Поэтому

\[
5t^3+7t^3
=
(5+7)t^3
=
12t^3.
\]

Выражения

\[
5t^3
\qquad\text{и}\qquad
7t^2
\]

подобными не являются.

Поэтому непосредственно сложить их нельзя:

\[
5t^3+7t^2.
\]

Можно вынести общий множитель:

\[
5t^3+7t^2
=
t^2(5t+7).
\]

При умножении действует уже правило степеней:

\[
5t^3\cdot7t^3
=
35t^6.
\]

\key{Не смешивайте два разных действия:}

\[
5t^3+7t^3=12t^3,
\]

\[
5t^3\cdot7t^3=35t^6.
\]

% =========================================================

\section*{13. Дополнительная задача по планиметрии}
\sectionline

Четырёхугольник \(ABCD\) вписан в окружность. Прямые \(AB\) и \(CD\) пересекаются в точке \(K\).

Дано:

\[
BK=6,
\qquad
DK=10,
\qquad
BC=12.
\]

Найти \(AD\).

\begin{center}
\begin{tikzpicture}[scale=0.45]
  \tkzDefPoint(0,0){K}
  \tkzDefPoint(6,0){B}
  \tkzDefPoint(13.3333,0){A}
  \tkzDefPoint(-3.6667,7.1096){C}
  \tkzDefPoint(-4.5833,8.8870){D}

  \tkzDefCircle[circum](A,B,C)
  \tkzGetPoint{O}

  \tkzDrawCircle[color=accent,line width=0.9pt](O,A)
  \tkzDrawSegments[color=textgray,line width=0.9pt](K,A K,D B,C A,D)
  \tkzDrawPoints(K,A,B,C,D)

  \tkzMarkAngle[
    size=0.8cm,
    color=accent,
    line width=0.8pt
  ](B,C,K)

  \tkzMarkAngle[
    size=0.8cm,
    color=accent,
    line width=0.8pt
  ](B,A,D)

  \tkzMarkAngle[
    size=0.65cm,
    color=darkaccent,
    line width=0.7pt
  ](B,K,C)

  \tkzLabelPoints[below](K,B,A)
  \tkzLabelPoints[left](C,D)

  \tkzLabelSegment[below](K,B){$6$}
  \tkzLabelSegment[left](K,D){$10$}
  \tkzLabelSegment[above left](B,C){$12$}
  \tkzLabelSegment[above right](A,D){$x$}
\end{tikzpicture}
\end{center}

Рассматриваем треугольники

\[
\triangle BCK
\qquad\text{и}\qquad
\triangle ADK.
\]

Имеем

\[
\angle BKC=\angle AKD,
\]

поскольку лучи \(KB\) и \(KA\), а также \(KC\) и \(KD\) попарно совпадают.

Кроме того,

\[
\angle BAD+\angle BCD=180^\circ
\]

по свойству вписанного четырёхугольника.

Углы \(BCK\) и \(BCD\) смежные, поэтому

\[
\angle BCK+\angle BCD=180^\circ.
\]

Следовательно,

\[
\angle BCK=\angle BAD.
\]

Получаем подобие по двум углам:

\[
\triangle BCK\sim\triangle ADK.
\]

Соответствующие стороны:

\[
BC\longleftrightarrow AD,
\qquad
BK\longleftrightarrow DK.
\]

Поэтому

\[
\frac{AD}{BC}
=
\frac{DK}{BK}.
\]

Подставляем значения:

\[
\frac{x}{12}
=
\frac{10}{6}.
\]

Отсюда

\[
6x=120,
\]

\[
\boxed{x=20}.
\]

\key{Полезная идея:}
если две секущие проходят через вписанный четырёхугольник и пересекаются вне окружности, ищите пару подобных треугольников.

% =========================================================

\section*{14. Типичные ошибки}
\sectionline

\begin{itemize}
  \item Приравнивать показатели слишком рано.
  \item Оставлять дополнительный множитель рядом со степенью и после этого «отбрасывать» основания.
  \item Забывать правило
  \[
  a^{-n}=\frac1{a^n}.
  \]
  \item Ошибаться при переводе корня:
  \[
  \sqrt[k]{a^m}=a^{m/k}.
  \]
  \item При возведении степени в степень складывать показатели вместо умножения.
  \item Применять формулы степеней к сумме или разности.
  \item Терять скобки в сложном показателе.
  \item Ошибаться со знаком при раскрытии скобок.
  \item Сокращать отдельные слагаемые алгебраической дроби.
  \item Не учитывать условие ненулевого знаменателя.
  \item Смешивать сложение подобных слагаемых и умножение степеней.
  \item Записывать показатель степени так низко или далеко от основания, что запись становится неоднозначной.
\end{itemize}

% =========================================================

\section*{15. Мини-чек перед завершением решения}
\sectionline

\begin{tabularx}{\linewidth}{@{}p{0.32\linewidth}X@{}}
\toprule
\textbf{Проверка} & \textbf{Контрольный вопрос} \\
\midrule
Основание &
Можно ли представить обе части через одно основание? \\

Составные числа &
Все ли числа вроде \(9\), \(27\), \(64\), \(25\) преобразованы удобным образом? \\

Корни &
Верно ли применена формула дробной степени? \\

Дроби &
Не требуется ли отрицательная степень? \\

Степень степени &
Перемножены ли показатели? \\

Структура &
Осталась ли слева и справа ровно одна степень перед приравниванием показателей? \\

Сумма или разность &
Можно ли вынести общий степенной множитель? \\

Алгебра &
Нет ли ошибки в знаках, скобках и переносах? \\

ОДЗ &
Есть ли знаменатель или другое выражение, создающее ограничение? \\

Оформление &
Чётко ли видны основание и показатель каждой степени? \\
\bottomrule
\end{tabularx}

% =========================================================
% ТРЕНИРОВОЧНЫЙ БЛОК
% =========================================================

\clearpage

\section*{Тренировочный блок}
\sectionline

\textbf{Формат: Путь Б.}

По 5 заданий на каждый день.

Решения в тренировочном блоке не приводятся.

% ---------------------------------------------------------

\subsection*{08.09.2026}

\begin{enumerate}[label=\textbf{\arabic*)}]

  \item Решите уравнение
  \[
  2^{3x-4}=32.
  \]

  \item Решите уравнение
  \[
  5^{\frac{x}{2}+1}=\frac1{25}.
  \]

  \item Решите уравнение
  \[
  3^{2x+1}=\sqrt{27}.
  \]

  \item Решите уравнение
  \[
  3\cdot9^{x-2}=\frac1{27}.
  \]

  \item Решите уравнение
  \[
  4^{x+1}-4^x=48.
  \]

\end{enumerate}

% ---------------------------------------------------------

\subsection*{09.09.2026}

\begin{enumerate}[label=\textbf{\arabic*)},start=6]

  \item Решите уравнение
  \[
  \left(\frac13\right)^{2x-5}
  =
  \frac1{81}.
  \]

  \item Решите уравнение
  \[
  \left(\frac12\right)^{4-x}=8.
  \]

  \item Решите уравнение
  \[
  \frac{5^{x+2}}{\sqrt{125}}
  =
  5^{3/2}.
  \]

  \item Решите уравнение
  \[
  2\cdot8^{x-1}
  =
  \frac1{16}.
  \]

  \item Решите уравнение
  \[
  9^{x+1}-9^x=648.
  \]

\end{enumerate}

\vspace{1em}

\textbf{Дополнительное закрепление ошибки из домашней работы}

Упростите выражение и укажите допустимые значения переменной:

\[
\frac{8y^{10}+12y^{10}}{12y^{10}}.
\]

\vspace{0.7em}

\textbf{Дополнительное геометрическое закрепление}

Четырёхугольник \(ABCD\) вписан в окружность, прямые \(AB\) и \(CD\) пересекаются в точке \(K\).

Известно:

\[
BK=8,
\qquad
DK=12,
\qquad
BC=14.
\]

Найдите \(AD\).

% =========================================================
% ОТВЕТЫ
% =========================================================

\clearpage

\section*{Ответы к тренировочному блоку}
\sectionline

\begin{table}[ht]
\centering
\caption{Ответы к заданиям 08.09--09.09.2026}

\begin{tabularx}{\linewidth}{@{}cX@{}}
\toprule
\textbf{№} & \textbf{Ответ} \\
\midrule

1 & \(x=3\) \\

2 & \(x=-6\) \\

3 & \(x=\dfrac14\) \\

4 & \(x=0\) \\

5 & \(x=2\) \\

6 & \(x=\dfrac92\) \\

7 & \(x=7\) \\

8 & \(x=1\) \\

9 & \(x=-\dfrac23\) \\

10 & \(x=2\) \\

\bottomrule
\end{tabularx}
\end{table}

\vspace{1em}

\begin{table}[ht]
\centering
\caption{Ответы к дополнительным заданиям}

\begin{tabularx}{\linewidth}{@{}p{0.34\linewidth}X@{}}
\toprule
\textbf{Задание} & \textbf{Ответ} \\
\midrule

Алгебраическая дробь &
\[
\frac53,
\qquad
y\ne0
\]
\\

Геометрия &
\[
AD=21
\]
\\

\bottomrule
\end{tabularx}
\end{table}

\end{document}